\section{Introdução}

A emissão de certificados digitais é uma atividade essencial para empresas de treinamentos, consultorias e instituições educacionais que promovem cursos, capacitações e eventos. Apesar da relevância desse processo, grande parte das organizações ainda depende de procedimentos manuais ou soluções improvisadas, utilizando ferramentas como Google Planilhas, Google Docs, PowerPoint e complementos externos.

Essas práticas resultam em retrabalho, inconsistências, falta de padronização e dificuldade de controle operacional. Além disso, ferramentas como AutoCrat e documentos espalhados entre diferentes plataformas não foram projetados para uso profissional em escala, o que gera limitação, fragilidade e risco de inconsistências.

A ausência de um sistema integrado compromete a rastreabilidade e a credibilidade dos certificados, uma vez que não há mecanismos oficiais de validação pública, histórico centralizado, indicadores gerenciais ou emissão em lote eficiente. Esse cenário se agrava com o crescimento de cursos online e da necessidade de certificação confiável.

Diante disso, justifica-se o desenvolvimento de uma solução centralizada e automatizada, capaz de reduzir erros humanos, otimizar o tempo operacional e garantir segurança, conformidade e profissionalismo. Recursos como códigos únicos, validação via QR Code e assinaturas digitais ampliam a confiabilidade dos certificados e fortalecem a imagem institucional das organizações emissoras.

A existência de uma empresa parceira piloto, que atualmente utiliza soluções limitadas, reforça a aplicabilidade prática deste projeto e evidencia a demanda real do mercado por sistemas robustos de emissão e autenticação de certificados digitais.


\subsection{Objetivos}

\subsubsection{Geral}

Desenvolver um sistema web automatizado para emissão, gestão, envio e validação de certificados, substituindo processos manuais e descentralizados por uma plataforma profissional, segura e escalável.

\subsubsection{Específicos}

\begin{itemize}
    \item automatizar a criação de certificados por meio de templates configuráveis;
    \item integrar o sistema com Google Sheets, Excel ou arquivos CSV para emissão em lote;
    \item possibilitar envio automático por e-mail com variáveis dinâmicas;
    \item implementar mecanismos de autenticação pública via código único e QR Code;
    \item disponibilizar dashboards e estatísticas de emissão;
    \item manter histórico centralizado e rastreável de todas as emissões.
\end{itemize}
